
\vspace{0.5cm}
\noindent
\begin{tabularx}{\textwidth}{|>{\raggedright\arraybackslash}p{4cm}|X|}
\hline
\textbf{Use-case ID} \newline \textbf{Use-case name} & U4.1 \newline \textbf{Cập nhật trạng thái vị trí đỗ} \\ \hline
\textbf{Use-case overview} & Cho phép hệ thống IoT thu thập và cập nhật liên tục tình trạng của từng vị trí đỗ xe lên hệ thống trung tâm. \\ \hline
\textbf{Actors} & 1. IoT Sensors/Gateway \\ \hline
\textbf{Preconditions} & 
1. Hệ thống đang hoạt động bình thường (không mất kết nối, không bị quá tải)\newline
2. Các cảm biến IoT tại vị trí đỗ đã được cấp nguồn và có kết nối mạng ổn định tới Gateway. \\ \hline
\textbf{Trigger} & Cứ mỗi 15s cảm biến thu thập và gửi thông tin về sự thay đổi vật lý tại vị trí đỗ. \\ \hline
\textbf{Steps} & 
1. Cảm biến quét khu vực đỗ và cập nhật trạng thái.\newline
2. IoT Gateway tiếp nhận tín hiệu từ cảm biến và gắn thêm mốc thời gian cụ thể.\newline
3. IoT Gateway truyền dữ liệu trạng thái cùng định danh của vị trí đỗ về hệ thống quản lý trung tâm.\newline
4. Hệ thống trung tâm tiếp nhận, lưu trữ dữ liệu vào cơ sở dữ liệu và xác nhận quá trình cập nhật hoàn tất. \\ \hline
\textbf{Post conditions} & Trạng thái mới nhất của vị trí đỗ xe được lưu trữ chính xác trong cơ sở dữ liệu của hệ thống. \\ \hline
\textbf{Exception flow} & 
- \textbf{E1 (Mất kết nối mạng):} IoT Gateway không thể gửi dữ liệu lên hệ thống. Dữ liệu trạng thái sẽ được lưu cục bộ tại Gateway và tự động đồng bộ lại ngay khi kết nối được khôi phục. \newline 
- \textbf{E2 (Dữ liệu không hợp lệ)}: Nếu hệ thống phát hiện gói tin sai cấu trúc hoặc cảm biến không tồn tại trong danh mục, hệ thống sẽ từ chối lưu bản ghi và tự động ghi lại lỗi để Admin kiểm tra. \\ \hline
\end{tabularx}
\newpage




\vspace{0.5cm}
\noindent
\begin{tabularx}{\textwidth}{|>{\raggedright\arraybackslash}p{4cm}|X|}
\hline
\textbf{Use-case ID} \newline \textbf{Use-case name} & U4.2 \newline  \textbf{Hiển thị trạng thái bãi xe} \\ \hline
\textbf{Use-case overview} & Cung cấp cho User cái nhìn tổng quan theo thời gian thực về tình trạng sức chứa của bãi xe (số chỗ trống, số chỗ đã đầy, sơ đồ vị trí). \\ \hline
\textbf{Actors} & 1. User \\ \hline
\textbf{Preconditions} & 
1. Hệ thống đang hoạt động bình thường (không mất kết nối, không bị quá tải).\newline
2. User đã đăng nhập thành công vào ứng dụng/tài khoản của bãi xe.\newline
3. Hệ thống đã thu thập được dữ liệu trạng thái từ các IoT Gateway.\\ \hline
\textbf{Trigger} & User truy cập vào chức năng "Trạng thái bãi xe". \\ \hline
\textbf{Steps} & 
1. Hệ thống tiếp nhận yêu cầu xem trạng thái bãi xe từ User.\newline
2. Hệ thống truy xuất cơ sở dữ liệu để lấy dữ liệu trạng thái mới nhất của tất cả các vị trí đỗ.\newline
3. Hệ thống hiển vị trí chính xác của những vị trí chỗ trống/đã có xe và xuất lên giao diện dưới dạng sơ đồ bãi đỗ, màu xanh tương đương với trống và màu đỗ là đã có xe. \\ \hline
\textbf{Post conditions} & Giao diện màn hình của User hiển thị chính xác và trực quan trạng thái hiện tại của toàn bộ bãi đỗ xe. \\ \hline
\textbf{Exception flow} & 
- \textbf{E1 (Lỗi kết nối với CSDL trung tâm):} Hệ thống không thể lấy dữ liệu, hiển thị thông báo "Không thể tải sơ đồ bãi xe lúc này, vui lòng thử lại sau ít phút".\newline
- \textbf{E2 (Dữ liệu không đồng bộ):} Nếu có độ trễ lớn từ IoT Gateway, hệ thống sẽ hiển thị cảnh báo "Dữ liệu có thể không phải là mới nhất". \\ \hline
\end{tabularx}
\newpage




\vspace{0.5cm}
\noindent
\begin{tabularx}{\textwidth}{|>{\raggedright\arraybackslash}p{4cm}|X|}
\hline
\textbf{Use-case ID} \newline \textbf{Use-case name} & U4.3 \newline \textbf{Thông báo cảm biến lỗi} \\ \hline
\textbf{Use-case overview} & Hệ thống tự động theo dõi tình trạng hoạt động của các thiết bị IoT và cảnh báo cho Admin khi phát hiện thiết bị có dấu hiệu hư hỏng hoặc mất kết nối. \\ \hline
\textbf{Actors} & IoT Sensors/Gateway, Admin. \\ \hline
\textbf{Preconditions} & 
1. Chức năng giám sát tình trạng thiết bị (Heartbeat/Ping) của hệ thống đang chạy nền. \newline
2.Toàn bộ thiết bị IoT đã được đăng ký mã định danh vào hệ thống.\\ \hline
\textbf{Trigger} & Hệ thống trung tâm không nhận được tín hiệu phản hồi từ cảm biến/Gateway trong khoảng thời gian cho phép (45s), hoặc nhận được mã lỗi từ thiết bị gửi về. \\ \hline
\textbf{Steps} & 
1. Hệ thống ghi nhận sự cố từ một thiết bị phần cứng cụ thể.\newline
2. Hệ thống cập nhật trạng thái của thiết bị đó thành "Error" trong cơ sở dữ liệu.\newline
3. Hệ thống tự động tạo một thông báo cảnh báo chi tiết (bao gồm mã thiết bị, vị trí, thời gian xảy ra sự cố).\newline
4. Hệ thống đẩy thông báo hiển thị lên màn hình của Admin. \newline
5. Admin xác nhận thông báo, cập nhật trạng thái "Đang bảo trì" cho thiết bị đó trong cơ sở dữ liệu.\\ \hline
\textbf{Post conditions} & Thông tin lỗi thiết bị được ghi vào hệ thống và Admin nhận được thông báo cảnh báo kịp thời. \\ \hline
\textbf{Exception flow} & 
- \textbf{E1:} Nếu Admin không thực hiện việc xác nhận, cảnh báo vẫn sẽ được lưu trữ và bôi đỏ trực tiếp trên giao diện Dashboard. \newline
- \textbf{E2 (Thiết bị tự phục hồi kết nối)}: Nếu thiết bị bị mất kết nối mạng tạm thời nhưng sau đó tự kết nối lại được trước khi Admin kịp xác nhận báo lỗi, hệ thống tự động thu hồi thông báo lỗi ban đầu và thay bằng một thông báo mới "Hệ thống IoT đã tự khôi phục kết nối thành công". \\ \hline
\end{tabularx}
\newpage
